\documentclass[11pt]{article}
\usepackage[T1]{fontenc}
\usepackage[utf8]{inputenc}
\usepackage{lmodern,microtype}
\usepackage[margin=1.08in]{geometry}
\usepackage{amsmath,amssymb,amsthm,mathtools}
\usepackage{xcolor,hyperref}
\hypersetup{colorlinks=true,linkcolor=blue!55!black,citecolor=blue!55!black,urlcolor=blue!60!black}
\setlength{\parindent}{0pt}
\setlength{\parskip}{0.48em}
\setlength{\emergencystretch}{2em}
\newtheorem{theorem}{Theorem}
\theoremstyle{remark}
\newtheorem{remark}[theorem]{Remark}
\newcommand{\spt}{\operatorname{spt}}

\title{Infinite-Order Contact Does Not Force Containment for Nonintegral Stationary Varifolds}
\author{[Author Name]}
\date{}

\begin{document}
\maketitle

\begin{abstract}
\textbf{Relation to the problem list.}
This manuscript concerns Question~30 (L30), ``Infinite-order contact forces
containment,'' but it is \emph{not} a solution of the intended
Brena--Decio--De~Lellis problem.  Their problem is for integer rectifiable
varifolds, whereas the question in the list omits integrality.  The example
below uses noninteger weights essentially.  It disproves only the broader
statement printed in the list; the intended integral-varifold problem remains
open.

We give an explicit stationary rectifiable varifold whose support has
infinite-order quadratic contact with a plane at the origin but is not
locally contained in that plane.  The construction is a base plane together
with a rapidly weighted sequence of parallel planes converging to it.  The
weights are nonintegral; this is essential to local finiteness.  The example
shows that an infinite-order-contact rigidity statement for stationary
varifolds requires an integrality or comparable multiplicity hypothesis.
\end{abstract}

\section{Introduction}

Let $V$ be an $m$-dimensional stationary rectifiable varifold in a Euclidean
open set and let $\mathcal M$ be a smooth minimal $m$-submanifold through
$x_0\in\spt V$.  A natural unique-continuation question asks whether
\begin{equation}\label{eq:contact}
 \int_{B_r(x_0)}\operatorname{dist}(x,\mathcal M)^2\,d\|V\|(x)
 =o(r^N)
 \qquad\text{for every }N\in\mathbb N
\end{equation}
forces $\spt V$ to be locally contained in $\mathcal M$.  Brena, Decio,
and De Lellis formulate this problem for integer rectifiable varifolds
\cite{BDD}.  We show that the corresponding assertion for arbitrary
real-weighted rectifiable varifolds is false.

\begin{theorem}\label{thm:main}
For every integer $m\ge1$, there exist a stationary rectifiable
$m$-varifold $V$ in $\mathbb R^{m+1}$, an $m$-plane $P$, and a point
$0\in P\cap\spt V$ such that \eqref{eq:contact} holds with
$\mathcal M=P$, but
\[
 \spt V\not\subset P
\]
in every neighborhood of the origin.  The varifold $V$ is locally Radon
but is not integral.
\end{theorem}

\section{Construction and proof}

Let
\[
 P=\mathbb R^m\times\{0\}\subset\mathbb R^{m+1}
\]
and fix a unit normal vector $e$.  For $j\ge1$, set
\[
 h_j=2^{-j},
 \qquad
 a_j=e^{-h_j^{-2}}=e^{-4^j},
 \qquad
 P_j=P+h_je.
\]
Define
\begin{equation}\label{eq:V}
 V=|P|+\sum_{j=1}^{\infty}a_j|P_j|.
\end{equation}

\begin{proof}[Proof of Theorem~\ref{thm:main}]
First we verify that \eqref{eq:V} defines a rectifiable Radon varifold.
For every compact set $K\subset\mathbb R^{m+1}$, the $m$-dimensional
measure of $K\cap P_j$ is bounded uniformly in $j$, while
$\sum_j a_j<\infty$.  Hence the series of weight measures converges on
compact sets.  Its support is contained in the countably rectifiable set
$P\cup\bigcup_jP_j$, and its multiplicity is measurable and positive.

Every affine plane is stationary.  If $X$ is a compactly supported smooth
vector field, then
\[
 \delta V(X)=\delta|P|(X)+\sum_{j=1}^{\infty}a_j\delta|P_j|(X)=0.
\]
The series is absolutely convergent because the masses of the planes on
$\spt X$ are uniformly bounded and $\sum_j a_j<\infty$.  Thus $V$ is
stationary.

Let $\omega_m$ denote the volume of the unit ball in $\mathbb R^m$, and
put
\[
 I(r)=\int_{B_r(0)}\operatorname{dist}(x,P)^2\,d\|V\|(x).
\]
The base plane contributes zero.  If $h_j<r$, then
$P_j\cap B_r(0)$ is an $m$-ball of radius
$\sqrt{r^2-h_j^2}$ and every point on it has distance $h_j$ from $P$.
Consequently
\begin{equation}\label{eq:I}
 I(r)=\omega_m\sum_{h_j<r}
 a_jh_j^2(r^2-h_j^2)^{m/2}.
\end{equation}

Let $J=J(r)$ be the least index satisfying $h_J<r$.  Since
$4^{j+1}\ge4^j+1$ for $j\ge1$, the superexponential tail obeys
\[
 \sum_{j\ge J}e^{-4^j}\le C e^{-4^J}
\]
for an absolute constant $C$.  Moreover $h_J<r$, so
$4^J=h_J^{-2}>r^{-2}$.  Using \eqref{eq:I}, we obtain
\[
 I(r)
 \le \omega_m r^{m+2}\sum_{j\ge J}a_j
 \le C_m r^{m+2}e^{-1/r^2}.
\]
Since $e^{-1/r^2}=o(r^k)$ as $r\downarrow0$ for every $k>0$, it follows
that $I(r)=o(r^N)$ for every $N\in\mathbb N$.

On the other hand, $h_j\downarrow0$.  Every ball centered at the origin
therefore meets some parallel plane $P_j$, and every such $P_j$ belongs to
$\spt V$.  Hence the support is not contained in $P$ in any neighborhood
of the origin.

Finally, the multiplicity of $V$ on $P_j$ is
$a_j=e^{-4^j}\notin\mathbb Z$.  Thus $V$ is not an integral varifold.
\end{proof}

\begin{remark}
The fractional weights are the mechanism that makes the example locally
finite.  Replacing them by positive integer multiplicities on infinitely
many full parallel planes would produce infinite mass in every fixed ball.
Accordingly, the construction does not resolve the integral-varifold unique
continuation problem in \cite{BDD}.
\end{remark}

\begin{thebibliography}{9}

\bibitem{BDD}
C.~Brena, S.~Decio, and C.~De Lellis,
\emph{Remarks and conjectures on stationary varifolds},
Rend. Lincei Mat. Appl. \textbf{36} (2025), 409--443;
\href{https://arxiv.org/abs/2503.00651}{arXiv:2503.00651}.

\bibitem{Simon}
L.~Simon,
\emph{Lectures on Geometric Measure Theory},
Proceedings of the Centre for Mathematical Analysis, vol.~3,
Australian National University, Canberra, 1983.

\end{thebibliography}
\end{document}
